“Hallo lieve Emmy,
na alle commotie,

aangaande ons vertrek
in zo’n kort tijdsbestek,

stuur ik je deze brief”,
schreef Emmy haar zuslief.

Het was de enige
brief van de menige,

die Emma beloofde
en Emmy geloofde.

Emmy schreef een paar keer,
maar kreeg geen briefverkeer.

“Ga je die verbranden?”
De trillende handen

van de snode vader,
brachten de hoop nader

tot zijn kloppende borst.
Desondanks nagevorst

zag Tante de naam Emmy.
De correspondentie

bleek volledig voor haar.
Tante begreep het gevaar.

Vader kon niet lezen
en alleen maar vrezen

dat de vele brieven
Emmy zouden grieven

en vader zijn daden
vast zouden verraden.

Alsof hij als Phoenix
door verbranden een fix

zocht om zo uit de as
als zondaar die hij was

te worden herboren.
Zijn list leek verloren.

Ook een tweede stapel
lag klaar op de kapel,

boven de openhaard.
Emma stond op elke kaart.

Emma beschreef de grond,
hoe de afmars ontstond

en die was verlopen.
Dat zij en ma hopen

haar spoedig weer te zien.
Want ze miste sindsdien

voortdurend haar zuslief.
“Ik stuur je deze brief

vanwege een voorstel.
Mama zit in de knel,

omdat onze landheer
wil gaan schenden haar eer

even als die van mij.
Aldus vertelde zij.

Ze laat papa schieten
en laat hem genieten,

zonder te weigeren,
laat staan te steigeren.

De deur moest dan op slot
voor elk fysiek genot

met mij als oudste kind,
opdat ik me verbind

en mijn hand wordt gevraagd.
Ze liegt: ‘Ik ben nog maagd.’

en kiest mijn echtgenoot.”
Dat de landheer besloot

om niet zijn herenrecht
te eisen voor de echt,

bleek verder uit de brief.
Ook moeders ongerief

over een valse list,
waar ze nadien van wist.

De dubbele binding,
die de landheer ontving,

veroorzaakte zijn nijd.
Haar tegenstrijdigheid

leverde verwarring
en wekte verstarring.

Met zijn lijfeigenschap
als geplande misstap

ontsloeg de landheer
haar vader van de eer

kosteloos het bordeel
grenzend aan zijn kasteel

te mogen gebruiken.
Vader wou ontduiken

de pijn door ontmoeting
van soms zijn gemalin

of zijn oudste dochter.
Vader kon en mocht er

hen niet tegenkomen.
Daardoor werd voorkomen

vader zijn invloedssfeer
door de sluwe landheer.

Het adres van zijn vrouw
en dochter koos hij flauw

gelijk aan het bordeel
en niet van zijn kasteel.

Dus brieven meegeven
verried het seksleven

dat vader eerder had.
Zo werd verhinderd dat

Emmy haar berichten
de junctuur ontwrichtten

naar de grap verluidde.
Michel Foucault duidde

het adres van de brief
als een dispositief

of machtsmechanisme
door pa’s conformisme.

“Een eega komt met spoed
uit een ander landgoed

om met mij te trouwen.
Ik moet maar vertrouwen

dat de landheer aldaar
mij niet voor het altaar

al neemt als profiteur
door zijn droit de seigneur.”

Emmy kende de plot,
waarin steeds het genot,

ieder vertrouwen schond.
Toen de schout voor haar stond

en vroeg waar vader was
gaf ze geen tegengas,

omdat ze geen gevaar
zag in haar handgebaar.

Tot ze va zag huilen
met meerdere builen
en haar fout besefte.
Hij die de pacht hefte

kwam de schuld verhalen.
Pa kon niet betalen.

De landheers soldaten
konden dat niet laten.

Haar vader werd geveld
door symbolisch geweld

zoals Bourdieu zei.
In elke maatschappij

zit een ongelijkheid
een harde gewoonheid

wiens zwakke consensus
hij zag als habitus.

De steeds stijgende pacht,
die de landheer bedacht,

moest worden verrekend.
Er al op gerekend

liet hij zich maar knechten
en dumpte zijn rechten.

Werd de onteigene
zo een lijfeigene.

Evenals zijn gezin.
“Ik en ma zijn slavin!”,

schreef Emma daarover.
En de goedgelover

nam daarom de benen.
Uit het zicht verdwenen

ontsnapte ze haar lot:
“Gevangen in genot",

beschreven als tanha
in de Dhammapada.

“De landheer ging testen
door mama te pesten

of ze haar belofte
niet nadien verslofte.

Ze moest hem masseren
om te profileren

hoe gedreven ze was
op het Eemse matras.

En niet ouwehoeren
of een act opvoeren.

Alles werd ruim verkend
zoals een happy end.

Ze deed het allemaal
bloot in de ridderzaal

vol met vele gasten,
die haar ook betastten.

Ze kende de meesten.
Het waren de beesten,

die haar ook aanzagen
bij de landheers plagen

tijdens het bedacht
recht van de eerste nacht."

Emmy vroeg zich toch af
of wat Emma aangaf

de ter plaatse adel
zag als acceptabel.

Het strookte met de praat,
vol met laster en smaad,

die als koppig gerucht
steeds rondzwierf in de lucht.

Er werd tersluiks geklaagd
dat toen ze werd ontmaagd

er omstanders waren.
Die zagen haar paren

op een weelderig bed,
waar spiegels zijn gezet.

In stoelen gezeten
moest het publiek weten

als luide getuige
hoe hij haar aftuigde.

Suf in hun dronken bui.
Want de beste stuurlui

stonden weer aan de wal.
Moeder was de speelbal

in meerdere keren
van machtige heren.

Was dit moreel beleid?
Was dit rechtvaardigheid?

Confucius gaf aan
dat elke onderdaan

moest kunnen rekenen
op schone tekenen

van deugdzaam leiderschap.
Roddel en achterklap,
onrust of een opstand
hield men zo in de hand.

Vader speelde een rol.
Hij zag die als eervol.

Hij hielp seksverslaafden,
die in lust door draafden

binnen zijn parochie.
Dat dacht althans Emmy,

omdat vader dat zei.
“Ja, die taak geloofd hij,

en wij laten hem gaan,
in die naïeve waan!”,

vertelde de pastoor.
“Je vader heeft niet door

dat hij zelf wordt gesteund,
daar hij zelf miskleunt.

Redden van anderen
kan hem veranderen!”

De wrange situatie
door landheers sommatie

van moeder en Emma
bleef zo een dilemma.

Zou Emmy ooit weten:
“Je bent niet vergeten”,

zoals Emma eens schreef
in een brief die wegbleef.
